\textbf{Classroom test, part 2.} You have 40 minutes in total.
Three answer-sheet tables are provided. Tables 2.1ab and 2.1c correspond to part 2.1, ``Messier Objects'',
and Table 2.2 corresponds to part 2.2, ``Bright Stars''. Mark answers in these tables and divide your time as you wish.
For both parts, only the first 10 answers in each column are evaluated if more are given.
There is no negative marking for wrong answers.

\textbf{Part 2.1, Messier Objects.}
Table 2.1ab lists all 110 Messier objects, with its first row filled in as an example.
\begin{description}
\item[(a)] Write the letter code of the correct type for any 10 objects in the corresponding ``Type'' column. \hfill\mbox{[5 points]}
Choose Galaxy (GL), Nebula (NB), Globular Cluster (GC), Open Cluster (OC), or other (OT).
A few Messier objects fall outside the first four categories and are classed as other.
\item[(b)] Write the IAU constellation code for any 10 objects, not necessarily the same ones as in (a), in the ``Constellation'' column. \hfill\mbox{[5 points]}
\item[(c)] Table 2.1c lists 25 objects also present in the Messier list.
For any 10, write their Messier number. \hfill\mbox{[10 points]}
\end{description}
\textbf{Part 2.2, Bright Stars.}
Table 2.2 lists the 92 brightest stars, excluding the Sun, up to 2.50 mag. It gives proper names, Bayer
designations, apparent visual magnitudes $m_v$, and distances from the Sun in light-years. The stars are ordered alphabetically.
\begin{description}
\item[(d)] Identify any 10 stars that are binary or members of multiple systems.
Tick the ``Binary or Multiple'' column. \hfill\mbox{[5 points]}
For systems, the given visual magnitude is the combined naked-eye magnitude.
\item[(e)] Identify any 10 variable stars, of any kind, by ticking the ``Variable'' column. \hfill\mbox{[10 points]}
The given magnitude for a variable is its approximate time-averaged naked-eye visual magnitude.
\item[(f)] The Morgan--Keenan (MK) spectral sequence uses O, B, A, F, G, K, M from hottest to coolest.
For this exercise G, K and M stars are called ``cool''. Identify any 10 cool stars and tick the ``Cool'' column. \hfill\mbox{[10 points]}
For systems in this list, either all members have the same type (for example all cool), or their brightest
component dominates the observed colour. A cool member may be much brighter than a hot member, giving a red or orange naked-eye appearance.
\end{description}
The blank Messier answer tables and the supplied bright-star data follow.
\ioaafig{2023_OBS-SKY_Q14-f1.pdf}
\ioaafig{2023_OBS-SKY_Q14-f2.pdf}
\ioaafig{2023_OBS-SKY_Q14-f3.pdf}
\ioaafig{2023_OBS-SKY_Q14-f4.pdf}
\ioaafig{2023_OBS-SKY_Q14-f5.pdf}
